\begin{frame}[t]{Основные результаты Sims--Zha}
\small

\vspace{0.3em}

\begin{columns}[T,onlytextwidth]

\begin{column}{0.31\textwidth}
\textbf{Импульсные отклики}

\vspace{0.3em}
\footnotesize
Идентифицированные монетарные шоки дают осмысленные реакции переменных, но эффект на выпуск ограничен.
\end{column}

\begin{column}{0.31\textwidth}
\textbf{Историческая декомпозиция}

\vspace{0.3em}
\footnotesize
Неожиданные сдвиги политики не объясняют основную часть исторических рецессий.
\end{column}

\begin{column}{0.31\textwidth}
\textbf{Контрфакты}

\vspace{0.3em}
\footnotesize
Правила с фиксированными $M$ или $R$ не дают существенного сглаживания делового цикла.
\end{column}

\end{columns}

\vspace{1.0em}

\begin{block}{Главный вывод}
\small
Рецессии в основном не являются следствием автономных монетарных шоков; важнее систематическая реакция политики на состояние экономики.
\end{block}

\end{frame}
